\section{Introduction}\label{sec:intro}

Reservoir computing trains only a linear readout: the input drives a fixed nonlinear dynamical system that unfolds the input
history into high-dimensional features~\cite{jaeger2001,jaeger2004}. Because the dynamical system itself is not trained, it is
well suited to physical hardware~\cite{tanaka2019}, and photonics is one of the most studied platforms~\cite{vandersande2017},
from delay systems with a single nonlinear node~\cite{appeltant2011} to passive silicon chips~\cite{vandoorne2014}.
Frequency-multiplexed schemes use the frequency components of light as nodes, obtaining tens of features from one optical path;
programmable spectral filters then set input and output weights per frequency
component~\cite{butschek2020,lupo2021,lupo2023}. Optical next-generation reservoir computing follows a related
idea~\cite{cox2024,wang2025ngrc}.

Kerr microresonators are natural candidates for such schemes. Four-wave mixing transfers pump power to tens of comb lines,
providing a strong on-chip nonlinearity, and the line spacing is wide enough to demultiplex and detect the lines one by
one~\cite{kippenberg2018,pasquazi2018}. Several Kerr-cavity reservoirs have been proposed recently, in microresonators and in fibre cavities (Table~\ref{tab:related}): a chaotic
comb whose pump power is modulated~\cite{shishavan2025}, a single soliton whose pump detuning is modulated~\cite{cuevas2025}, a
cavity soliton whose pump phase is modulated~\cite{arabieh2026soliton}, and a frequency-multiplexed fibre cavity driven by a
synchronous pulse train~\cite{arabieh2026pulse}. In all of them the signal enters through a single degree of freedom of the drive.
In the chaotic regime the comb lines grow out of noise through modulation instability (MI), so the same input can give different
outputs in two runs, which violates the consistency a reservoir needs~\cite{uchida2004,jaeger2001}
(Section~\ref{sec:mechanism}). In the soliton regime the output is reproducible, but the device must stay inside the soliton
existence range~\cite{herr2014}, which limits how strongly the input may perturb it.

\begin{table}[htbp]
\centering
\caption{Kerr-cavity and frequency-multiplexed reservoirs related to this work. ``Input'' is the degree of freedom carrying the signal.}
\label{tab:related}
\footnotesize
\begin{tabular}{p{3.3cm}p{2.9cm}p{2.6cm}p{1.6cm}p{3.6cm}}
\toprule
work & input & comb state & experiment & relation to this work \\
\midrule
Shaabani Shishavan et al.~\cite{shishavan2025} & pump power & chaotic (MI) & no & reimplemented as ``power'' encoding \\
Cuevas et al.~\cite{cuevas2025} & pump detuning & single soliton & yes & reimplemented as ``detuning'' encoding \\
Arabieh et al.~\cite{arabieh2026soliton} & pump phase & cavity soliton & yes & reimplemented as ``phase'' encoding \\
Arabieh et al.~\cite{arabieh2026pulse} & synchronous pulse train & pulse-driven cavity & no & closest to our uniform comb \\
Butschek, Lupo et al.~\cite{butschek2020,lupo2021,lupo2023} & spectral weights on comb lines & fibre loop, no Kerr comb & yes & per-line weighting, not in a Kerr resonator \\
this work & random complex weight per input line & coherently seeded Kerr comb & no & -- \\
\bottomrule
\end{tabular}
\end{table}

Per-line spectral weighting is therefore not new in itself~\cite{butschek2020,lupo2021}, and driving Kerr cavities with pulse
trains or combs is known~\cite{obrzud2017}; to our knowledge, driving a Kerr resonator with many random-phase coherent input lines
that sit on its resonances has not been studied as a reservoir computer. A second gap is the comparison with software. The works
above differ in how operating points are selected, on which noise realisations results are reported and how much noise the software
baseline receives, and we find that the last of these alone can reverse the conclusion of the comparison
(Section~\ref{sec:software}). It therefore remains unclear \emph{how well a Kerr-microcomb reservoir performs under realistic noise, what limits it, and how it compares with
software when the comparison is made under one protocol}.

We study these questions with simulations. Our contributions are:
\begin{enumerate}
\item \emph{A common benchmark} (Sections~\ref{sec:methods} and~\ref{sec:encodings}): eight encodings compared with one simulator,
      one set of device and power ranges, one readout, one noise model and one operating-point budget, with points chosen and
      reported on independent noise realisations.
\item \emph{A spectral-phase input mask for Kerr resonators} (Sections~\ref{sec:mechanism} and~\ref{sec:encodings}): random complex
      weights on 16 input lines give coherently seeded, reproducible dynamics and Mackey--Glass errors 15 to 45 times below a
      uniformly modulated comb and two to three orders of magnitude below our implementations of single-pump encodings; the randomness of the line phases matters more than per-line
      amplitudes. On channel equalisation the mask is comparable to a uniformly modulated comb, and with a single resonator it does
      not help NARMA-10. A linear-cavity control ($\gamma=0$) attributes the Mackey--Glass gain mostly to the Kerr dynamics;
      channel equalisation does not need them.
\item \emph{The noise floor and the laser requirement} (Sections~\ref{sec:noise} and~\ref{sec:readout}): with ideal detectors about
      half of the floor is detuning jitter, which a linear readout given the exact jitter sequence did not remove, so the laser
      linewidth must be a few kilohertz or less for our device; readout-side changes lower the error by about a factor of two.
\item \emph{A software calibration} (Section~\ref{sec:software}): the best resonator found matches a 31-node echo state network (ESN)
      fed through a converter with about 7 effective bits, and software stays ahead when the resonator's readout is quantized too.
\item \emph{Two directions and an optimised mask} (Section~\ref{sec:bidir}): the signal and the jitter appear to act through the
      same variable, the effective detuning, and none of the operating points we examined avoids the jitter; driving the two
      directions of one ring with inverted signals, reading out both (62 lines) and optimising the mask lowers the error 21-fold
      relative to the best single-direction ring (4.5-fold relative to one direction read out and optimised in the same way), to
      about 8 effective bits against an equally optimised ESN, which narrows but does not close the gap.
\end{enumerate}
Section~\ref{sec:gradient} reports a negative result on gradient-based design, and Section~\ref{sec:discussion} discusses an
experimental implementation, limitations and unmodelled effects.
